% !TeX program = lualatex
% Review-master workspace for literature mapping and later Introduction drafting.
% Keep this file long and source-oriented; move only checked points into
% manuscript-ready notes.
% Build in VS Code with latexmk + lualatex + BibTeX.

\documentclass[11pt]{article}

\usepackage{amsmath}
\usepackage{iftex}
\usepackage[defaultsups]{newtxtext}
\usepackage[varg]{newtxmath}
\usepackage[superscript]{cite}
\ifLuaTeX
\usepackage[luatex,pdfencoding=auto]{hyperref}
\else
\usepackage[dvipdfmx]{hyperref}
\usepackage{pxjahyper}
\fi

\hypersetup{hidelinks}

\title{Review Master}
\author{}
\date{\today}

\begin{document}

\maketitle
\tableofcontents

This document is a working review master for literature mapping, field understanding, and later Introduction drafting. Keep source-oriented notes here first, and move only checked points into the manuscript-ready notes at the end of each section.

\section{Pumped-storage and upper reservoir operation}
% Track why upper-reservoir hydraulics matter, what operating constraints shape
% the problem, and how reversible operation is framed in prior work.

\subsection{Question-oriented notes}
Use this subsection to group notes by question so that later Introduction paragraphs can be drafted from themes rather than from individual papers.
% Organize this section around questions that may later become Introduction
% paragraphs.
%
% Suggested block:
% - question:
% - current consensus:
% - disagreement or unresolved point:
% - papers to revisit:

\subsection{Cross-paper synthesis}
Record consensus, disagreement, and the remaining gap across multiple papers before deciding what belongs in the manuscript.
% Compare common assumptions, differences in reservoir geometry or operating
% context, and any gap between hydropower operation studies and basin-scale flow
% observations.

\subsection{Paper notes}
Use this subsection for source-by-source notes after checking the original paper directly.
% Duplicate the block below for each paper that belongs in this section.
%
% \subsubsection*{Short paper label}
% bibkey:
% status: candidate / checked / used
% question:
% method:
% main findings:
% relevance to this paper:
% caveat:
%
% After the point is verified, restate it in prose and cite it with \cite{...}.

\subsection{Manuscript-ready notes}
Distill only checked points that are worth carrying into the eventual Introduction.
% What from this section should move into the present paper's Introduction?
% Keep only points that are checked and useful.
%
% Suggested block:
% - working claim:
% - supporting citations:
% - why it matters for this paper:
% - what to leave out from the final Introduction:

\section{Shallow basin / reservoir flow patterns}
% Track reported basin-scale circulation patterns, recirculation cells, stagnant
% zones, depth dependence, and regime-transition language used in related work.

\subsection{Question-oriented notes}
Use this subsection to group notes by question so that later Introduction paragraphs can be drafted from themes rather than from individual papers.
% Use this subsection to group papers by hydraulic question rather than by paper.
%
% Suggested question types:
% - What flow patterns recur in shallow basins or shallow reservoirs?
% - How do geometry, depth, or discharge change those patterns?
% - Which patterns are steady, transitional, or intermittent?

\subsection{Cross-paper synthesis}
Record consensus, disagreement, and the remaining gap across multiple papers before deciding what belongs in the manuscript.
% Summarize common patterns, contradictory observations, and missing comparisons
% that matter for the present shallow single-port basin problem.

\subsection{Paper notes}
Use this subsection for source-by-source notes after checking the original paper directly.
% Duplicate the block below for each paper that belongs in this section.
%
% \subsubsection*{Short paper label}
% bibkey:
% status: candidate / checked / used
% question:
% method:
% main findings:
% relevance to this paper:
% caveat:

\subsection{Manuscript-ready notes}
Distill only checked points that are worth carrying into the eventual Introduction.
% Distill the few points that should later support the Introduction narrative on
% shallow-basin flow regimes and why a regime-oriented description is needed.

\section{Reversible inflow-outflow and mode asymmetry}
% Track studies that distinguish inflow from outflow, reversible use of the same
% port or structure, and mechanisms that create directional asymmetry.

\subsection{Question-oriented notes}
Use this subsection to group notes by question so that later Introduction paragraphs can be drafted from themes rather than from individual papers.
% Suggested question types:
% - When do inflow and outflow produce qualitatively different basin responses?
% - What mechanisms are proposed for asymmetry under the same geometry?
% - Which comparisons are direct, and which are only inferred?

\subsection{Cross-paper synthesis}
Record consensus, disagreement, and the remaining gap across multiple papers before deciding what belongs in the manuscript.
% Summarize where the literature already treats mode asymmetry explicitly and
% where the present study may fill a comparison gap.

\subsection{Paper notes}
Use this subsection for source-by-source notes after checking the original paper directly.
% Duplicate the block below for each paper that belongs in this section.
%
% \subsubsection*{Short paper label}
% bibkey:
% status: candidate / checked / used
% question:
% method:
% main findings:
% relevance to this paper:
% caveat:

\subsection{Manuscript-ready notes}
Distill only checked points that are worth carrying into the eventual Introduction.
% Draft concise claims that may later become the asymmetry part of the
% Introduction. Keep them short enough that they can be compressed further.

\section{Surface PIV and physical modeling}
% Track how surface-flow measurements are obtained, what physical-model choices
% are reported, and what quality limits are stated for surface PIV.

\subsection{Question-oriented notes}
Use this subsection to group notes by question so that later Introduction paragraphs can be drafted from themes rather than from individual papers.
% Suggested question types:
% - What can surface PIV resolve reliably in shallow free-surface experiments?
% - How are rectification, frame spacing, and surface tracers handled?
% - What uncertainty or interpretation limits are emphasized?

\subsection{Cross-paper synthesis}
Record consensus, disagreement, and the remaining gap across multiple papers before deciding what belongs in the manuscript.
% Compare measurement setups, preprocessing choices, and reporting practices.
% Note where methods are transferable to this project and where they are not.

\subsection{Paper notes}
Use this subsection for source-by-source notes after checking the original paper directly.
% Duplicate the block below for each paper that belongs in this section.
%
% \subsubsection*{Short paper label}
% bibkey:
% status: candidate / checked / used
% question:
% method:
% main findings:
% relevance to this paper:
% caveat:

\subsection{Manuscript-ready notes}
Distill only checked points that are worth carrying into the eventual Introduction.
% Keep only the measurement points that are worth carrying into the final
% Introduction or brief Methods motivation.

\section{Scaling and similitude}
% Track which similarity arguments are used, which scale effects are acknowledged,
% and how strongly each paper claims prototype relevance.

\subsection{Question-oriented notes}
Use this subsection to group notes by question so that later Introduction paragraphs can be drafted from themes rather than from individual papers.
% Suggested question types:
% - Which nondimensional controls are treated as primary?
% - Where do authors explicitly relax strict similitude?
% - What limitations are reported for shallow free-surface scaling?

\subsection{Cross-paper synthesis}
Record consensus, disagreement, and the remaining gap across multiple papers before deciding what belongs in the manuscript.
% Compare stronger and weaker similitude claims, note disagreements, and record
% language that can help frame this study as prototype-informed rather than as a
% strict full-similitude model.

\subsection{Paper notes}
Use this subsection for source-by-source notes after checking the original paper directly.
% Duplicate the block below for each paper that belongs in this section.
%
% \subsubsection*{Short paper label}
% bibkey:
% status: candidate / checked / used
% question:
% method:
% main findings:
% relevance to this paper:
% caveat:

\subsection{Manuscript-ready notes}
Distill only checked points that are worth carrying into the eventual Introduction.
% Write short, careful notes on how much scaling justification the final paper
% actually needs and what cautionary wording should remain.

\section{Sedimentation / CFD linkage}
% Track how flow-regime observations are linked to sedimentation or CFD, and keep
% clear boundaries between present evidence and future-study implications.

\subsection{Question-oriented notes}
Use this subsection to group notes by question so that later Introduction paragraphs can be drafted from themes rather than from individual papers.
% Suggested question types:
% - How do prior studies connect surface-flow patterns to sediment behavior?
% - How are laboratory observations used to motivate or validate CFD?
% - Which claims are direct results and which are only hypotheses?

\subsection{Cross-paper synthesis}
Record consensus, disagreement, and the remaining gap across multiple papers before deciding what belongs in the manuscript.
% Record transferable links, weak links, and future-facing gaps. Keep this
% subsection especially strict about status so that unchecked ideas do not slip
% into the manuscript narrative.

\subsection{Paper notes}
Use this subsection for source-by-source notes after checking the original paper directly.
% Duplicate the block below for each paper that belongs in this section.
%
% \subsubsection*{Short paper label}
% bibkey:
% status: candidate / checked / used
% question:
% method:
% main findings:
% relevance to this paper:
% caveat:

\subsection{Manuscript-ready notes}
Distill only checked points that are worth carrying into the eventual Introduction.
% Note only the minimal bridge needed for the present paper. Separate what is
% supported now from what belongs to future sedimentation or CFD studies.

\bibliographystyle{unsrt}
\bibliography{review_refs}

% Usage:
% 1. Add only verified bibliography entries to review_refs.bib.
% 2. Use status = candidate until the original paper is checked directly.
% 3. Move a note to status = used only after deciding it belongs in the paper.
% 4. Keep paper-specific detail in "Paper notes" and cross-paper logic above it.
% 5. Distill only checked points into each "Manuscript-ready notes" subsection.
% 6. When drafting the actual Introduction, copy from manuscript-ready notes
%    first, then shorten and merge across sections.
% 7. If a claim is still indirect or uncertain, leave it out of manuscript prose.

\end{document}
